\documentclass[10pt,a4paper]{article}
\usepackage[utf8]{inputenc} % para poder usar tildes en archivos UTF-8
\usepackage[T1]{fontenc}
\usepackage[spanish]{babel} % para que comandos como \today den el resultado en castellano
\addto\captionsspanish{\renewcommand{\tablename}{Tabla}} % tabla en vez de cuadro
\usepackage{a4wide} % márgenes un poco más anchos que lo usual
\usepackage[sinEntregas]{caratula}
\usepackage{booktabs}
\usepackage{float}
\usepackage{xcolor}
\usepackage{listings}
\usepackage[scaled=0.75]{roboto-mono}
\usepackage[pdfview=Fit]{hyperref} % Siempre preferentemente al final de los paquetes
\usepackage[all]{hypcap}
% Definición del fondo gris claro
\definecolor{codegray}{rgb}{0.96,0.96,0.96}

% Cambiamos 'Listing' por 'Código'
\renewcommand{\lstlistingname}{Código}

% Configuración del estilo para Python
\lstdefinestyle{colabstyle}{
    language=Python,
    basicstyle=\ttfamily\small,
    keywordstyle=\color{blue}\bfseries,
    commentstyle=\color{gray}\itshape,
    stringstyle=\color{purple},
    backgroundcolor=\color{codegray},  % Fondo gris claro
    frame=single,                       % Marco alrededor del bloque
    rulecolor=\color{black!15},         % Color del borde
    framesep=6pt,                      % Relleno (padding) interno para que el código no toque los bordes
    keepspaces=true,             
    showstringspaces=false,       
    captionpos=b,                   
    breaklines=true,
    literate={á}{{\'a}}1 {é}{{\'e}}1 {í}{{\'i}}1 {ó}{{\'o}}1 {ú}{{\'u}}1 {ñ}{{\~n}}1
}

% Aplicamos el estilo como predeterminado para todos los listings
\lstset{style=colabstyle}

\begin{document}

\titulo{Trabajo Práctico}

\fecha{\today}

\materia{Aprendizaje Automático I}
\grupo{Grupo 11}

\integrante{Bravo, Iván Agustín}{786/23}{bravoivanagustin@gmail.com}
\integrante{Kiss, María Delfina}{72/23}{kissdelfina@gmail.com}
\integrante{Lützow Holm, Eric}{323/12}{elholm90@gmail.com}
\integrante{Mendoza, Juan Cruz}{004/01}{email4@dominio.com}
\integrante{Rostan, Marcos Andrés}{589/23}{rostanmarcos@gmail.com}
% Pongan cuantos integrantes quieran

\maketitle
\section{Introducción}
En este trabajo práctico evaluamos distintos modelos de aprendizaje automático para determinar si una muestra de células tiene pronóstico de crecimiento anómalo o descontrolado, condición conocida como hiperplasia. El dataset consta de 500 instancias provenientes de pacientes con lesiones pretumorales, en las que se midió la expresión de 200 genes. Cada instancia está etiquetada como verdadero (buen pronóstico) si no hubo indicios de una nueva hiperplasia, o falso (mal pronóstico) en caso contrario.

Probamos diferentes algoritmos con distintas configuraciones de hiperparámetros, con el fin de comparar su comportamiento en esta tarea de clasificación: árboles de decisión, K vecinos más cercanos (KNN), support vector machines (SVM), análisis discriminante lineal (LDA) y Naive Bayes. La implementación se realizó en Python.\footnote{El notebook con el código completo se puede consultar \href{https://github.com/kissdelfina/aprendizaje-automatico/blob/main/Trabajo\%20Practico\%20I/TP_Aprendizaje_Automatico_I.ipynb}{aquí}. En este informe incluimos solo el código mínimo necesario para ilustrar los pasos más importantes.}

\section{Separación de datos}
Para estimar la performance del modelo final dividimos el dataset en dos particiones: un conjunto de desarrollo, sobre el cual entrenamos y seleccionamos los modelos, y un conjunto de control, para evaluar al final del trabajo la configuración elegida. Siguiendo el estándar habitual, destinamos el 15\% de las instancias al conjunto de control.

La partición debe ser representativa del dataset, la proporción de etiquetas tiene que coincidir con la del conjunto de desarrollo, y el conjunto de control no debería presentar propiedades ausentes en el resto de los datos, ni a la inversa.

Del total, 359 instancias (72\%) corresponden a la etiqueta 0 (mal pronóstico) y 141 (28\%) a la etiqueta 1 (buen pronóstico). Para preservar esas proporciones aplicamos muestreo estratificado: agrupamos las instancias por etiqueta y extrajimos una muestra aleatoria del 15\% dentro de cada grupo. Lo implementamos con los métodos \texttt{groupby} y \texttt{sample} de Pandas, fijando una semilla para que la partición sea reproducible. El resultado son $0{,}15 \cdot 359 = 54$ instancias de etiqueta 0 y $0{,}15 \cdot 141 \approx 21$ de etiqueta 1, lo que deja un conjunto de control de 75 instancias y un conjunto de desarrollo de 425.
 
Con la separación establecida, construimos los modelos.

\section{Construcción de modelos}
Buscamos construir y evaluar modelos de árboles de decisión, obteniendo una estimación realista de su performance. 
Como primer paso separamos el conjunto de desarrollo en entrenamiento y validación de manera proporcional a sus etiquetas, usando \texttt{train\_test\_split} de \texttt{sklearn} estratificado, y nos quedamos con un 0,8 de entrenamiento y 0,2 de validación.
Sobre esa partición entrenamos un árbol con altura máxima 3 y el resto de los hiperparámetros en sus valores por defecto.

De este árbol obtuvimos un accuracy de 0.6471 sobre el subconjunto de validación.
Ese valor sale de una única partición de 85 instancias, por lo que depende bastante de cuáles hayan caído del lado de validación.
Para obtener una estimación más estable repetimos la medición con validación cruzada sobre el conjunto de desarrollo completo, con un K-fold estratificado de 5 particiones de manera que cada fold conserve la proporción de etiquetas.
Para cada fold calculamos Accuracy, Area Under the Precision-Recall Curve (AUPRC) y Area Under the Receiver Operating Characteristic Curve (AUCROC), y además el score global utilizando únicamente los folds de validación.
Los resultados se pueden observar en la tabla \ref{tab:resultados_permutaciones}.

\begin{table}[H]
\centering
\resizebox{\textwidth}{!}{%
\begin{tabular}{c|cc|cc|cc}
\toprule
\textbf{Permutación} & \textbf{Acc (train)} & \textbf{Acc (val)} & 
\textbf{AUPRC (train)} & \textbf{AUPRC (val)} & 
\textbf{AUCROC (train)} & \textbf{AUCROC (val)} \\
\midrule
1 & 0.8265 & 0.6000 & 0.7153 & 0.3265 & 0.8626 & 0.5697 \\
2 & 0.8265 & 0.7412 & 0.6862 & 0.4213 & 0.8037 & 0.6411 \\
3 & 0.8353 & 0.6235 & 0.6981 & 0.2895 & 0.8506 & 0.4754 \\
4 & 0.8265 & 0.6471 & 0.6900 & 0.3594 & 0.8429 & 0.6117 \\
5 & 0.7971 & 0.6824 & 0.6491 & 0.3574 & 0.8538 & 0.5683 \\
\midrule
\textbf{Promedios} & \textbf{0.8224} & \textbf{0.6588} & 
\textbf{0.6878} & \textbf{0.3508} & 
\textbf{0.8427} & \textbf{0.5732} \\
\midrule
\textbf{Global} & \textbf{(NO)} & \textbf{0.6588} & 
\textbf{(NO)} & \textbf{0.3404} & 
\textbf{(NO)} & \textbf{0.5610} \\
\bottomrule
\end{tabular}%
}
\caption{Métricas de un árbol de altura máxima 3 con criterio Gini, por fold de la validación cruzada estratificada de 5 particiones sobre el conjunto de desarrollo. Global se calcula sobre las predicciones acumuladas de los 5 folds de validación}
\label{tab:resultados_permutaciones}
\end{table}

Para interpretar estos valores necesitamos un punto de comparación, así que evaluamos con la misma validación cruzada un clasificador que ignora los atributos y predice siempre la clase mayoritaria. Alcanza un accuracy de 0.7176, un AUPRC de 0.2824 que no es otra cosa que la proporción de positivos, y un AUCROC de 0.5000 porque asigna el mismo score a todas las instancias.

Luego exploramos 6 combinaciones de altura máxima y criterio de corte, manteniendo validación cruzada con K = 5 y generando las combinaciones con \texttt{ParameterGrid} de scikit-learn. 
Probamos alturas máximas de 3, 5 e infinita, con criterio de corte Gini y Entropía.
Reutilizamos las mismas particiones para que los valores sean comparables, lo que se verifica en que la fila de altura 3 con Gini reproduce exactamente el resultado de la tabla \ref{tab:resultados_permutaciones}. 
Los resultados se muestran en la tabla \ref{tab:accuracy_arboles}

\begin{table}[H]
\centering
\begin{tabular}{llcc}
\toprule
\textbf{Altura máxima} & \textbf{Criterio de corte} & \textbf{Accuracy (training)} & \textbf{Accuracy (validación)} \\
\midrule
3        & Gini     & 0.8224 & 0.6588 \\
5        & Gini     & 0.9482 & 0.6447 \\
Infinito & Gini     & 1.0000 & 0.6400 \\
3        & Entropía & 0.8029 & 0.6541 \\
5        & Entropía & 0.9253 & 0.6612 \\
Infinito & Entropía & 1.0000 & 0.5953 \\
\bottomrule
\end{tabular}
\caption{Accuracy promedio en entrenamiento y validacion de árboles de decisión, según altura máxima y criterio de corte, con la misma validación cruzada estratificada de 5 particiones sobre el conjunto de desarrollo.}
\label{tab:accuracy_arboles}
\end{table}

De la tabla \ref{tab:resultados_permutaciones} vemos que el modelo base, con altura máxima 3 y criterio Gini, logró un accuracy de validación de 0.6588, por debajo del 0.7176 del clasificador trivial.
Su AUCROC de validación es 0.5732 en promedio y 0.5610 global, apenas por encima del 0.5000 de referencia, y su AUPRC de 0.3508 mejora muy poco el 0.2824 que ya da predecir siempre la clase mayoritaria.
El modelo base entonces casi no captura señal, aunque la brecha con entrenamiento es amplia, 0.8224 contra 0.6588 en accuracy y 0.8427 contra 0.5732 en AUCROC.

En la tabla \ref{tab:accuracy_arboles} se observa el comportamiento típico del overfitting.
Al relajar la altura máxima el accuracy de entrenamiento sube hasta 1.0000, es decir que el árbol memoriza las 340 instancias de cada fold, mientras que el de validación no mejora.
Con Gini la caída es pareja, de 0.6588 a 0.6400, y con Entropía el peor caso es el árbol sin restricción, que baja a 0.5953.

Ninguna de las 6 configuraciones supera el 0.7176 del clasificador trivial, y el rango completo de la columna de validación, de 0.5953 a 0.6612, es más angosto que la variación entre folds de una misma configuración, que en la tabla \ref{tab:resultados_permutaciones} va de 0.6000 a 0.7412.
Por eso no consideramos que ninguna combinación sea preferible a las demás.
Tampoco el accuracy es la métrica más adecuada para este problema, porque con un 72\% de instancias negativas premia al modelo que ignora la clase minoritaria, que es justamente la que interesa detectar, y por eso de acá en adelante usamos AUCROC.
En conjunto, las dos tablas muestran que un árbol de decisión no alcanza para estos 200 atributos, lo que motiva explorar otras familias de algoritmos.

\section{Comparación de algoritmos}
En esta sección probamos árboles de decisión, KNN y SVM, con varias combinaciones de hiperparámetros. Usamos la AUCROC como métrica de comparación.

\subsection{Árboles de decisión}
Como hiperparámetros tuvimos en cuenta la profundidad máxima, la mínima cantidad de instancias que necesita un nodo para dividirse, la mínima cantidad de instancias de cada hoja, el criterio (Gini o Entropía) y la máxima cantidad de atributos, probando la raíz cuadrada y el logaritmo de la cantidad total.

\begin{lstlisting}
   parametros_arbol = {"max_depth": [3, 5, 10, 20, None],
                     "min_samples_split": [2, 5, 10, 20], 
                     "min_samples_leaf": [1, 5, 10, 15], 
                     "criterion": ["gini", "entropy"], 
                     "max_features": ["sqrt", "log2", None]} 

   arbol_base = DecisionTreeClassifier(random_state=random_state)

   buscador_arbol = RandomizedSearchCV(estimator=arbol_base, param_distributions=parametros_arbol,
                                       n_iter=50, scoring="roc_auc", cv=skf, random_state=random_state)

   buscador_arbol.fit(X_dev, y_dev)
\end{lstlisting}


Los mejores hiperparámetros encontrados fueron: 
\texttt{'min\_samples\_split': 10, 
'min\_samples\_leaf': 15, 
'max\_features': 'log2', 
'max\_depth': 5, 
'criterion': 'gini'}
Con un AUCROC score de 0.6740.

\subsection{KNN}
Para KNN, primero hay que tener en cuenta que suele ser mejor escalar las features. 
Como pueden tener escalas muy distintas, las features de mayor rango pueden llegar a dominar la distancia. Hiperparámetros:
\begin{itemize}
    \item \texttt{n\_neighbours}: mientras un k chico se puede llegar a ajustar al ruido local de cada datapoint, un k grande simplemente se aproxima a predecir la clase mayoritaria.

    \item \texttt{weights}: permite que los puntos más cercanos pesen más.

    \item \texttt{metric}: que es la distancia que usamos para medir cercanía.
\end{itemize}

\begin{lstlisting}
   parametros_knn = {"knn__n_neighbors": [3, 5, 7, 10, 15, 21, 31, 41, 51, 61, 75],
                     "knn__weights": ["uniform", "distance"], 
                     "knn__metric": ["euclidean", "manhattan"], 
                     "scaler": [StandardScaler(), None]}  
\end{lstlisting}

Mejores hiperparámetros encontrados para KNN: \texttt{'scaler': None, 'knn\_weights': 'distance', 'knn\_n\_neighbors': 21, 'knn\_metric': 'manhattan'}
El mejor score AUCROC fue 0.8474.

\subsection{SVM}
Para SVM, también necesitamos escalar los features. Hiperparámetros:

\begin{itemize}
    \item \texttt{C}: controla la regularización del SVM. Cuanto mayor es, obliga a SVM a tener un margen estricto, lo que puede llevar a overfitting, mientras que cuanto más chico es, tolera más puntos mal clasificados.
    \item \texttt{kernel}: nos dice si la frontera es lineal en el espacio de features, o se hace un mapping mediante el kernel trick para hacer un hiperplano en otro espacio.
    \item \texttt{degree}: para kernels polinómicos solamente debemos decidir hasta qué grado puede ser el polinomio.
    \item \texttt{gamma}: para el kernel RBF y poly, gamma define el radio de influencia de cada support vector.
\end{itemize}

Para \texttt{Gamma} y \texttt{C}, se suelen probar diferentes órdenes de magnitud. Para \texttt{degree}, grados mayores a 3 puede generar un espacio de features enorme y sobreajustar con pocas instancias.

\begin{lstlisting}
   parametros_svm = [
      {"svm__kernel": ["linear"], "svm__C": [0.01, 0.1, 1, 10, 100, 1000], "scaler": [StandardScaler(), None]},
      {"svm__kernel": ["rbf"], "svm__C": [0.01, 0.1, 1, 10, 100, 1000], "svm__gamma": ["scale", 0.0001, 0.001, 0.01, 0.1, 1], "scaler": [StandardScaler(), None]},
      {"svm__kernel": ["poly"], "svm__C": [0.01, 0.1, 1, 10, 100, 1000], "svm__degree": [2, 3], "svm__gamma": ["scale", 0.0001, 0.001, 0.01, 0.1, 1], "scaler": [StandardScaler(), None]}]
\end{lstlisting}

Los mejores hiperparámetros encontrados para SVM fueron \texttt{'svm\_kernel': 'poly', 'svm\_gamma': 1, 'svm\_degree': 2, 'svm\_C': 1000, 'scaler': None}
El mejor score AUCROC fue 0.8587.

Vemos que las top 5 mejores configuraciones tienen la misma AUCROC y sólo cambian el C y Gamma, indicando que sin escalar los mismos no tienen relevancia.

\begin{table}[H]
\centering
\resizebox{\textwidth}{!}{%
\begin{tabular}{c|c|c|c}
\toprule
\textbf{Índice} & \textbf{Params} & \textbf{Mean Test Score} & \textbf{Std Test Score} \\
\midrule
2  & \texttt{\{'svm\_\_kernel': 'poly', 'svm\_\_gamma': 1, ...\}} 
   & 0.858743 & 0.055193 \\
9  & \texttt{\{'svm\_\_kernel': 'poly', 'svm\_\_gamma': 0.0001, ...\}} 
   & 0.858743 & 0.055193 \\
23 & \texttt{\{'svm\_\_kernel': 'poly', 'svm\_\_gamma': 1, ...\}} 
   & 0.858743 & 0.055193 \\
29 & \texttt{\{'svm\_\_kernel': 'poly', 'svm\_\_gamma': 0.01, ...\}} 
   & 0.858743 & 0.055193 \\
25 & \texttt{\{'svm\_\_kernel': 'poly', 'svm\_\_gamma': 0.01, ...\}} 
   & 0.858743 & 0.055193 \\
\bottomrule
\end{tabular}%
}
\caption{Mejores configuraciones obtenidas mediante búsqueda de hiperparámetros.}
\label{tab:svm_gridsearch}
\end{table}

Resulta muy llamativo ver que tanto para SVM y KNN, las mejores configuraciones de hiperparámetros según AUCROC fueron aquellas que no utilizaron el Scaler. 
Esto indica que las features con un rango más grande terminaron siendo aquellas que también tenían mayor poder predictivo.

\subsection{LDA y Naive Bayes}
Por último, añadimos a la comparación LDA y Naive Bayes.

\begin{table}[H]
\centering
\begin{tabular}{c|cc}
\toprule
\textbf{Fold} & \textbf{LDA} & \textbf{Naive Bayes} \\
\midrule
1 & 0.7131 & 0.7425 \\
2 & 0.6919 & 0.7582 \\
3 & 0.6291 & 0.7589 \\
4 & 0.6086 & 0.6926 \\
5 & 0.5608 & 0.7357 \\
\midrule
\textbf{Promedio} & \textbf{0.6407} & \textbf{0.7376} \\
\bottomrule
\end{tabular}
\caption{AUCROC por fold y promedio para LDA y Naive Bayes.}
\label{tab:aucroc_lda_nb}
\end{table}

Vemos que el AUCROC de LDA es el peor con AUC = 0.6407. 
Esto puede deberse a que se está estimando una matriz de covarianza de 200x200 (por las 200 features con que se entrenan) con pocas instancias de train por fold, lo que resulta en una estimación ruidosa y casi singular. 
En este caso, podríamos utilizar el hiperparámetro de \texttt{shrinkage} que regulariza la matriz de covarianza.

Por otro lado, Naive Bayes, incluso sin haber hecho una búsqueda de hiperparámetros, tuvo un AUC = 0.7376, similar a la performance de SVM y KNN. 
Este modelo, al asumir independencia entre los features, sólo estima la media y varianza por cada gen, por lo que cada una de las estimaciones puede ser buena incluso con pocos datos. 
Uno de los hiperparámetros que podríamos modificar es \texttt{var\_smoothing}, el cual añade una fracción de la varianza más grande de las features a cada una de las estimaciones de varianza. 
Esto funciona como una regularización de la varianza, y también ayuda a que si hay una feature con varianza casi cero, un valor distinto en test no reciba probabilidad casi cero, dominando el producto.

\section{Evaluación de performance}

El mejor modelo que obtuvimos durante el desarrollo fue el entrenado con el algoritmo \textbf{COMPLETAR} con hiperparámetros \textbf{COMPLETAR}.

Para aproximar el desempeño de este modelo sobre el conjunto de control (en posesión de los profesores), entrenamos el modelo con todas las instancias del conjunto de desarrollo y luego lo evaluamos sobre el conjunto de evaluación que separamos al principio del trabajo. De esta forma obtuvimos una estimación del desempeño final del modelo en datos no vistos previamente, sin sobreestimarlo.

El valor de AUC-ROC obtenido por el modelo final sobre este conjunto fue de \textbf{COMPLETAR}. Los scores de cada instancia del conjunto de control se encuentran adjuntos con el trabajo, en el archivo \texttt{11\_y\_pred\_control\_\textbf{COMPLETAR}.csv}.

\section{Conclusiones}

\section{Discusión}

\section{Bibliografía consultada}

\end{document}